% Architecture générale.
%
% Principes de tracé, tenus strictement :
%   - flux principal vertical, une seule colonne : la requête se lit de haut
%     en bas sans jamais revenir en arrière ;
%   - la chaîne hors ligne occupe une colonne parallèle et ne rejoint la
%     colonne principale qu'en un seul point — l'étape de recherche, qui est
%     précisément le seul endroit où les deux pipelines se rencontrent ;
%   - routage orthogonal exclusivement, aucun segment oblique ;
%   - aucun croisement d'arêtes ;
%   - gabarits de nœuds uniformes par colonne.
\begin{tikzpicture}[
    font=\sffamily\scriptsize,
    boite/.style={draw, line width=0.7pt, rounded corners=2pt,
                  text width=34mm, minimum height=9.5mm, align=center,
                  inner sep=3pt},
    principal/.style={boite, minimum height=10.5mm},
    horsligne/.style={boite, text width=30mm, minimum height=8.5mm},
    d/.style={draw=slate,     fill=mist,          text=ink},          % donnée
    t/.style={draw=inptblue,  fill=inptblue!8,    text=navy},         % traitement
    m/.style={draw=leafgreen, fill=leafgreen!8,   text=leafgreen!40!black}, % modèle
    stop/.style={draw=anrtred, fill=anrtred!6,    text=anrtred!70!black},
    choix/.style={draw=amberdark, fill=amberdark!10, text=amberdark!55!black,
                  line width=0.7pt, diamond, aspect=2.3, align=center,
                  inner sep=1pt, text width=19mm},
    fleche/.style={-{Stealth[length=2.3mm,width=1.7mm]}, line width=0.7pt,
                   draw=slate, rounded corners=2.5pt},
    remise/.style={fleche, draw=inptblue!75, line width=0.9pt},
    refus/.style={fleche, draw=anrtred},
    tag/.style={font=\sffamily\scriptsize, text=slate, inner sep=2pt},
  ]

% ============================================================ grille
\def\XA{0}      % colonne principale (requête)
\def\XB{5.6}    % colonne hors ligne (indexation)
\def\XR{-4.5}   % dérivation d'abstention

% ============================================ colonne principale : la requête
\node[principal, d] (q)      at (\XA, 0)      {Question du citoyen};
\node[principal, m] (vec)    at (\XA,-1.75)   {Vectorisation\\BGE-M3};
\node[principal, t] (rdense) at (\XA,-3.5)    {Recherche dense};
\node[principal, t] (rlex)   at (\XA,-4.75)   {Recherche lexicale};
\node[principal, t] (rrf)    at (\XA,-6.4)    {Fusion RRF};
\node[choix]        (gate)   at (\XA,-8.05)   {score $\geq$\\\SeuilAbstention{} ?};
\node[principal, m] (gen)    at (\XA,-9.9)    {Génération\\Mistral 7B via Ollama};
\node[principal, t] (verif)  at (\XA,-11.65)  {Contrôle des citations};
\node[principal, d] (rep)    at (\XA,-13.4)   {Réponse citée\\+ avertissement};

\draw[fleche] (q)      -- (vec);
\draw[fleche] (vec)    -- (rdense);
\draw[fleche] (rlex)   -- (rrf);
\draw[fleche] (rrf)    -- (gate);
\draw[fleche] (gate)   -- node[tag, right=1pt, pos=0.35] {oui} (gen);
\draw[fleche] (gen)    -- (verif);
\draw[fleche] (verif)  -- (rep);

% Les deux recherches sont menées en parallèle sur la même question : la
% dérivation part du même point et les rejoint toutes deux.
\draw[fleche] (vec.west) -- ++(-1.05,0) |- (rlex.west);

% =================================== colonne hors ligne : la chaîne d'indexation
% Les deux index sont posés exactement sur les ordonnées des deux recherches
% qu'ils alimentent : la remise se fait par un segment horizontal pur.
\node[horsligne, d] (pdf)   at (\XB, 1.9)   {PDF officiel\\Loi 65-99};
\node[horsligne, t] (pars)  at (\XB, 0.5)   {\module{parser.py}\\extraction, découpage};
\node[horsligne, d] (jsonl) at (\XB,-0.9)   {Corpus JSONL\\\NbArticles{} articles};
\node[horsligne, m] (emb)   at (\XB,-2.3)   {Vectorisation\\BGE-M3};
\node[horsligne, d] (ichr)  at (\XB,-3.5)   {Index dense\\Chroma};
\node[horsligne, d] (ibm)   at (\XB,-4.75)  {Index lexical\\BM25};

\draw[fleche] (pdf)   -- (pars);
\draw[fleche] (pars)  -- (jsonl);
\draw[fleche] (jsonl) -- (emb);
\draw[fleche] (emb)   -- (ichr);
% BM25 indexe le texte, pas les vecteurs : il dérive du corpus, pas du modèle.
\draw[fleche] (jsonl.east) -- ++(1.15,0) |- (ibm.east);

% ============================ point de rencontre unique entre les deux chaînes
\draw[remise] (ichr.west) -- (rdense.east);
\draw[remise] (ibm.west)  -- (rlex.east);

% ================================================ dérivation d'abstention
\node[principal, stop] (abst) at (\XR,-8.05) {Abstention immédiate\\le modèle n'est pas appelé};
\draw[refus] (gate.west) -- node[tag, above, text=anrtred] {non} (abst.east);

% Aucune annotation en prose dans la figure : le commentaire appartient à la
% légende et au corps du texte, pas au schéma.

% Cadres de regroupement, posés en arrière-plan.
\begin{scope}[on background layer]
  \node[draw=rule, line width=0.7pt, rounded corners=3pt, dash pattern=on 3pt off 2pt,
        inner xsep=4mm, inner ysep=5mm, fit=(pdf)(pars)(jsonl)(emb)(ichr)(ibm)] (bhl) {};
  % L'abstention fait partie du pipeline de requête : le cadre l'englobe.
  \node[draw=rule, line width=0.7pt, rounded corners=3pt, dash pattern=on 3pt off 2pt,
        inner xsep=4mm, inner ysep=5mm,
        fit=(q)(vec)(rdense)(rlex)(rrf)(gate)(gen)(verif)(rep)(abst)] (ben) {};
\end{scope}

\node[font=\sffamily\scriptsize\bfseries, text=slate, anchor=south, align=center]
  at (bhl.north) {INDEXATION\\[-1pt]\scriptsize\mdseries hors ligne, une seule fois};
% Le cadre de requête étant élargi par la dérivation d'abstention, son titre
% se cale sur la colonne principale plutôt que sur le milieu du cadre.
\node[font=\sffamily\scriptsize\bfseries, text=slate, anchor=south, align=center]
  at (q.north |- ben.north) {REQUÊTE\\[-1pt]\scriptsize\mdseries en ligne, à chaque question};

\end{tikzpicture}
